\section{Methodology}

\subsection{Data Preparation}



CUDAHeat accepts a NetCDF source path or an \code{xarray.Dataset} containing the selected SST variable. The temporal coordinate is resolved explicitly and validated as a contiguous hourly sequence. All non-temporal dimensions are treated as spatial dimensions and are stacked into a single point axis for computation. This transformation is representational rather than analytical: point coordinates and the original grid indexing are retained so that event products can be returned in their native spatial geometry.
For path-based inputs, the package can construct or reuse a validated compact time-by-point NetCDF representation. The cache is held outside the source file and is identified by source and storage metadata, thereby preventing reuse after incompatible source changes. This optimization improves the regularity of spatial reads without modifying the scientific time axis. Inputs already arranged in time-by-point form are used directly; alternatively, users may bypass or require optimization explicitly.

\subsection{Hourly Climatology and Threshold Estimation}



Let $S(t,p)$ denote SST at time $t$ and point $p$. A baseline interval selects the rows used to estimate the reference cycle. Under the default \code{month_day_hour} calendar, each timestamp is assigned to one of $366\times24$ bins, preserving calendar date and clock hour. A \code{dayofyear} mode is also available. The default construction follows an hourly extension of the seasonal-threshold approach: for each calendar day and clock hour, samples from the same clock hour within $\pm5$ calendar days of the target day are pooled across the baseline. The climatological mean $C$ and percentile threshold $Q_q$ are calculated from this pool, with $q=0.90$ by default \citep{Hobday2016}.
For the month--day--hour calendar, the initially estimated 29 February bins are interpolated from the surrounding calendar days. This preserves a complete 366-day cycle while avoiding a discontinuity associated with the sparse occurrence of leap-day observations. All statistics are NaN-aware: a bin with no finite support remains missing rather than receiving an imputed temperature.



The raw mean and threshold cycles are circularly smoothed with a 31-day window by default, equivalent to 744 h. Smoothing is applied separately to each clock-hour sequence; it therefore regularizes the seasonal evolution of, for example, 14:00 SST without blending it with the 02:00 climatology. The smoothed cycles are then indexed by the full timestamp sequence to obtain $C(t,p)$ and $Q_q(t,p)$. The corresponding anomaly and exceedance magnitude are $A(t,p)=S(t,p)-C(t,p)$ and $E(t,p)=S(t,p)-Q_q(t,p)$, respectively.

\subsection{Event Detection}



The instantaneous exceedance mask is defined as
\[
M(t,p)=\bigl[S(t,p)>Q_q(t,p)\bigr]\land\bigl[S(t,p)\ \mathrm{is\ finite}\bigr].
\]
The strict inequality is intentional: observations equal to the threshold are not exceedances. A missing SST or threshold cannot generate an exceedance. CUDAHeat retains this raw mask conceptually distinct from the continuity-adjusted mask used to identify candidate intervals, allowing enrichment rules to measure raw thermal behaviour even when event continuity permits short interruptions.



The detector first fills only internal false runs bounded on both sides by exceedances and no longer than \code{max_internal_gap_hours}; the default is 6 h. Maximal true intervals are then extracted from this filled mask and retained when their duration is at least \code{min_duration_hours}, which defaults to 120 h. Finally, neighbouring duration-qualified intervals at the same spatial point are joined when their separating gap does not exceed \code{join_gap_hours}, which defaults to 48 h. The order of these operations is consequential: duration filtering precedes joining, and the joining operation acts on sparse intervals rather than on a dense label field.

The following parameter table summarizes the user-modulable settings that define the hourly reference and the temporal event rule. These values should be treated as an explicit analytical configuration, not as universally transferable constants.

\begin{table*}[t]
\centering
\footnotesize
\renewcommand{\arraystretch}{1.18}
\rowcolors{2}{gray!12}{white}
\begin{tabularx}{\textwidth}{@{}
    p{0.1\textwidth}
    p{0.10\textwidth}
    p{0.2\textwidth}
    X
@{}}
\toprule
Parameter & Default & Permitted setting & Scientific role \\
\midrule
\code{baseline} & \code{None} & \code{None} or \code{(start_year, end_year)} & Selects the observations used to estimate the climatological mean and percentile threshold. \\
\code{percentile} & $0.90$ & $0<q<1$ & Defines the quantile threshold $Q_q(t,p)$. \\
\code{calendar_mode} & \code{month_day_hour} & \code{month_day_hour} or \code{dayofyear} & Determines the assignment of timestamps to hourly climatology bins. \\
\code{smooth_window_hours} & 744 h & Positive 24-h multiple within the annual cycle & Sets the same-clock-hour circular smoothing window. \\
\code{smooth_min_periods} & 1 day & Integer from 1 to the smoothing-window length in days & Sets the minimum finite support required during smoothing. \\
\code{min_duration_hours} & 120 h & Integer $\geq 1$ & Retains candidate intervals at or above the prescribed duration. \\
\code{max_internal_gap_hours} & 6 h & Integer $\geq 0$ & Fills bounded sub-threshold interruptions before duration filtering. \\
\code{join_gap_hours} & 48 h & Integer $\geq 0$ & Joins neighbouring duration-qualified intervals at the same point. \\
\bottomrule
\end{tabularx}
\rowcolors{2}{}{}
\par\vspace{3pt}
\parbox{\textwidth}{\sffamily\small\textbf{Parameter settings.}\enspace User-modulable CUDAHeat parameters governing climatology construction and event detection; values shown are the current defaults.}
\end{table*}

\subsection{Event Representation}



Each retained event is represented as a point index and a half-open interval $[t_{\mathrm{start}},t_{\mathrm{stop}})$, together with duration and associated metadata. This sparse representation is appropriate because marine heatwave occupancy is ordinarily much smaller than the full time--space domain. It also supports direct event filtering by time, location or duration without reconstructing a dense cube.



When a dense representation is required for interoperability, CUDAHeat reconstructs a Dask-backed event-ID field in time chunks. Zero denotes the absence of a retained event and positive integers identify catalogue entries. The reconstruction is deliberately lazy: users can subset or write the result without allocating the full time--space array in host memory. Point-level and event-level accessors additionally provide SST, climatology, threshold and event context for targeted inspection.
